\section{\textcolor{red}{Целевая сегментация пользователей на орбитали}}
\label{sec:segmentation}

В настоящей главе описывается процедура сегментации пользователей на
орбитали, лежащая в основе предлагаемого подхода к оценке долгосрочного
эффекта. Основная идея состоит в том, чтобы использовать ансамбль
решающих деревьев, обученный на целевой метрике, и затем трансформировать
его в компактную и интерпретируемую систему правил, задающую устойчивые
поведенческие сегменты.

Под орбиталью далее понимается сегмент пользователей, определяемый
набором интерпретируемых правил в пространстве признаков и отражающий
характерный \textcolor{red}{архетип (режим)} поведения. Такая сегментация служит базовым
представлением пользовательской \textcolor{red}{базы} при анализе долгосрочного
эффекта через переходы между орбиталями.

\subsection{Постановка задачи сегментации}
\label{subsec:segmentation_setup}

\subsubsection{Общая схема сегментации}
\label{subsubsec:segmentation_pipeline}

Сегментация пользователей на орбитали строится \textcolor{red}{как последовательность
нескольких этапов}. Сначала выбирается целевая \textcolor{red}{метрика (величина)}, относительно
которой сегментация должна быть информативной. Затем для каждого
пользователя формируется набор признаков, описывающих его поведение
на фиксированном историческом интервале.

На следующем шаге обучается ансамбль решающих деревьев, предсказывающий
выбранную целевую метрику по указанным признакам. После обучения ансамбль
переводится в полиномиальное представление, из которого с помощью
фреймворка \texttt{Monoforest}~\cite{monoforest} извлекаются наиболее
значимые логические правила. На основе отобранных правил строится
сокращённая система сегментов, которой далее придаётся семантическая
интерпретация в виде орбиталей.

Таким образом, предсказательная модель используется не как конечный
инструмент оценивания, а как источник структуры, по которой строится
интерпретируемая сегментация пользователей.

\subsubsection{Определение целевой метрики}
\label{subsubsec:target_definition}

На первом этапе определяется целевая \textcolor{red}{метрика (величина)}, относительно которой
должна быть информативна сегментация. В качестве такой \textcolor{red}{метрики (величины)}
используется \textcolor{red}{прокси-метрика, связанная с будущей ценностью пользователя.
В рассматриваемой постановке в качестве примеров можно указать
ограниченную накопленную ценность пользователя или число активных дней
за фиксированный период времени.}
% Комментарий: формулировка нормальная, но можно чуть привязать к предыдущей главе,
% например "определённой в разделе ... ограниченной LTV" — сейчас связь имплицитна.

После выбора целевой метрики обучается ансамбль деревьев,
предсказывающий её по историческим признакам пользователя. Однако
в дальнейшем для построения сегментации используются не сами индивидуальные
предсказания модели, а наиболее существенные элементы её внутренней
структуры, которые извлекаются в виде логических правил.

Такой подход ориентирует сегментацию не на абстрактное сходство
пользователей, а на различия между ними с точки зрения будущей
целевой \textcolor{red}{метрики (величины)}.

\subsubsection{Формирование поведенческих признаков}
\label{subsubsec:feature_engineering}

Для каждого пользователя формируется набор признаков, описывающих его
поведение на фиксированном историческом интервале. В базовой
постановке используются признаки, рассчитанные за период длиной
30 дней.
% Комментарий: если в экспериментальной главе фактически другое окно (14/60 дней),
% лучше тут либо сослаться на конкретный выбор, либо сказать "например, 30 дней".

В качестве примеров поведенческих признаков можно привести:
\begin{itemize}
    \item \texttt{recency} — число дней с момента последней активности;
    \item \texttt{frequency} — число активных дней за наблюдаемый период;
    \item \texttt{GMV} — суммарную денежную выручку за период;
    \item \textcolor{red}{ и другие} агрегированные характеристики поведения пользователя.
\end{itemize}

Набор признаков проектируется отдельно для каждой системы с учётом
продуктовой специфики. Важно, чтобы он обеспечивал достаточно полное
описание среднесрочного поведения пользователя и позволял предсказывать
выбранную целевую метрику.

\subsection{Извлечение правил из ансамбля деревьев}
\label{subsec:rule_based_representation}

\subsubsection{Обучение ансамбля деревьев}
\label{subsubsec:model_training}

На сформированных признаках обучается ансамбль решающих деревьев,
реализуемый, в частности, в виде модели \texttt{CatBoostRegressor},
предсказывающей выбранную целевую \textcolor{red}{метрику (величину)}. \textcolor{red}{В качестве иллюстрации можно
рассматривать задачу предсказания числа активных дней пользователя на
следующем 30-дневном интервале.}

Использование ансамбля деревьев позволяет учитывать сложные
нелинейные зависимости между краткосрочным поведением и будущими
характеристиками пользователя. За счёт этого модель отражает
неоднородность пользовательской базы и выявляет структурные различия
между типами пользователей.

 \textcolor{red}{
Предполагается, что фундаментальные правила сегментации определяются
свойствами системы и остаются относительно стабильными во времени.}

% Комментарий: это сильное предположение; если останется в таком виде,
% хорошо бы в экспериментальной главе хоть как-то его обсудить/проверить.

\textcolor{red}{
Поэтому при обучении ансамбля допустимо использовать данные за
несколько временных периодов. Существенные изменения продукта, напротив,
требуют переобучения ансамбля и, как следствие, пересмотра набора правил.}


\subsubsection{Полиномиальное представление ансамбля деревьев}
\label{subsubsec:polynomial_representation}

Ключевой элемент фреймворка \texttt{Monoforest} заключается в том, что
ансамбль решающих деревьев может быть представлен не только как набор
деревьев, но и в виде полинома по индикаторам пороговых условий
на признаках~\cite{monoforest}. Такое представление позволяет перейти
от графовой структуры деревьев к единому набору интерпретируемых
логических правил.

Рассмотрим одно дерево решений \(h(x)\). Оно разбивает пространство
признаков на непересекающиеся области, соответствующие листьям дерева.
Если каждому листу \(l\) сопоставлено значение \(w_l\), то дерево можно
записать в виде
\begin{equation}
    h(x) = \sum_{l \in \mathrm{leaves}} w_l I\{x \in l\}.
\end{equation}

Индикатор принадлежности листу определяется набором условий на пути
от корня дерева к этому листу. Пусть элементарное условие имеет вид
\[
c(x) = I\{x_{i(c)} > b(c)\}.
\]
Тогда прохождению по правой ветви соответствует множитель \(c(x)\),
по левой — множитель \(1 - c(x)\). Индикатор каждого листа можно
представить в виде произведения таких множителей.

После раскрытия скобок дерево допускает представление в виде суммы
мономов:
\begin{equation}
    h(x) = \sum_{M \in 2^C} w_M \prod_{c \in M} c(x),
\end{equation}
где \(C\) — множество элементарных условий, а \(M\) — некоторое их
подмножество. Таким образом, каждое дерево представляется как сумма
элементарных логических правил, каждое из которых соответствует
комбинации пороговых условий.

Для ансамбля деревьев
\[
H(x) = \sum_{t=1}^{T} h_t(x)
\]
получается аналогичное полиномиальное представление:
\begin{equation}
    H(x) = \sum_{M \in 2^C} w_M \prod_{c \in M} c(x),
\end{equation}
где коэффициенты \(w_M\) агрегируют вклад одинаковых мономов,
возникающих в разных деревьях ансамбля~\cite{monoforest}.

Такое представление позволяет рассматривать ансамбль как единый
набор логических правил и служит основой для количественной оценки
значимости каждого правила и последующего отбора наиболее информативных.

\subsubsection{Извлечение и отбор логических правил}
\label{subsubsec:rule_extraction}

После перехода к полиномиальному представлению отдельные мономы могут
рассматриваться как интерпретируемые логические правила. Каждое такое
правило соответствует комбинации пороговых условий на признаках
пользователя и задаёт определённую область в пространстве признаков.

Формально моном записывается в виде
\begin{equation}
    M_k = \prod_j I\{x_{i(j)} \ge \theta_j\}.
\end{equation}
Произведение индикаторов означает одновременное выполнение всех
соответствующих условий, то есть моном задаёт правило вида: пользователь
относится к заданной области, если одновременно выполняется набор
пороговых ограничений на его признаки.

Например, правило может иметь вид
\[
I\{\texttt{frequency} \ge 10\} \cdot
I\{\texttt{recency} \le 3\} \cdot
I\{\texttt{GMV} \ge 5000\},
\]
что соответствует пользователям с высокой частотой активности, недавней
последней активностью и значительным уровнем монетизации.

Для оценки значимости правила \texttt{Monoforest} сопоставляет каждому мономому
вес
\begin{equation}
    \texttt{score}(M_k) = w_k^2 \cdot \mathrm{coverage}(M_k),
\end{equation}
где \(w_k\) характеризует вклад монома в значение модели, а
\(\mathrm{coverage}(M_k)\) — долю объектов выборки, для которых
выполняется соответствующее правило~\cite{monoforest}.

Такое определение веса сочетает две естественные характеристики
правила: величину влияния на предсказание и распространённость на
данных. Наибольший интерес представляют мономы с большими значениями
\(w_k\), которые одновременно существенно влияют на модель и выполняются
для заметной доли пользователей.

Все мономы ранжируются по величине весов \(\texttt{score}(M_k)\), после чего выбирается
ограниченный набор наиболее значимых правил. В результате полная
структура ансамбля заменяется компактным представлением через небольшое
число информативных логических конструкций.

\subsection{Построение и интерпретация орбиталей}
\label{subsec:orbital_interpretation}

\subsubsection{\textcolor{red}{Построение орбиталей по набору правил}}
\label{subsubsec:orbital_derivation}

После ранжирования мономов по весам \(\texttt{score}(M_k)\) выбирается верхняя часть
списка, например \(K\) наиболее значимых правил. Этот набор образует
сокращённое описание ансамбля, сохраняющее его ключевые структурные
элементы.

\textcolor{red}{
Далее по выбранным правилам строится упрощённая сегментационная
структура. Практически она может быть реализована в виде
необязательно бинарного дерева решений или эквивалентной системы
вложенных логических условий. Каждому конечному листу такой структуры
соответствует кандидатный сегмент пользователей.
}

Пользователь относится к сегменту в зависимости от того, какие из
отобранных правил выполняются для его признаков. Тем самым исходный
ансамбль решающих деревьев заменяется компактной системой сегментов,
определяемых небольшим набором наиболее значимых логических правил.

\subsubsection{\textcolor{red}{Семантическая разметка орбиталей}}
\label{subsubsec:semantic_labeling}

\textcolor{red}{
Полученные на предыдущем этапе сегменты целиком определяются данными
и структурой ансамбля, однако сами по себе не всегда имеют очевидный
содержательный смысл. Для придания сегментам интерпретируемости им
присваиваются осмысленные семантические метки, которые затем
используются в последующем анализе.}

В качестве примеров таких меток можно привести
\texttt{weak}, \texttt{active}, \texttt{paid},
\texttt{active\&paid} и \texttt{whale}. Отнесение сегмента к той или
иной категории основано на анализе распределений ключевых метрик
внутри сегмента.

Так, сегмент с высокой частотой активности, недавней последней
активностью и очень большим значением GMV естественно интерпретировать
как \texttt{whale}. Сегмент с умеренной активностью и положительной
монетизацией — как \texttt{active\&paid}. Сегменты с минимальной
активностью относятся к категории \texttt{weak}.

\textcolor{red}{
Отдельно выделяется сегмент \texttt{out}, включающий пользователей,
лежащих вне рассматриваемого исторического окна (например, не
проявлявших активности в этот период). Такой сегмент представляет
понятный, но вспомогательный поведенческий архетип, поэтому в дальнейшем
он используется главным образом для описания состава выборки и далее
не включается в интерпретацию основных орбиталей.
}

Полученное множество семантически размеченных сегментов образует
набор орбиталей, используемый в формальной модели и в
экспериментальном анализе.

\subsubsection{Свойства и роль орбиталей}
\label{subsubsec:orbital_properties}

Построенная сегментация обладает рядом свойств, важных для оценки
долгосрочного эффекта.

Во‑первых, она ориентирована на целевую метрику: орбитали формируются
как сегменты, информативные относительно будущей ценности пользователя,
а не как произвольные кластеры поведения.

Во‑вторых, орбитали интерпретируемы, поскольку задаются ограниченным
набором правил и снабжены содержательными метками. Это делает их
пригодными не только для вычислительных процедур, но и для анализа
со стороны продуктовых и бизнес‑команд.

В‑третьих, итоговая сегментация детерминирована и может применяться
к пользователям в режиме, близком к реальному времени.
Благодаря этому орбитали могут быть интегрированы в инфраструктуру
онлайн‑экспериментов.

Тем самым ансамбль деревьев, обладающий свойствами «чёрного ящика»,
преобразуется в прозрачную схему сегментации, основанную на
интерпретируемых правилах. Полученные орбитали служат базовой
единицей анализа, через которую далее описываются изменения
структуры пользовательской базы и оценивается долгосрочный эффект
экспериментального воздействия.
% Комментарий: последняя фраза чуть перекликается с концовкой формальной главы.
% Не ошибка, но можно сократить до одной фразы, если будет ощущение повтора.